\section{Discussion and Limitations}
\label{sec:discussion}

% SKELETON — write after final numbers land (Week 4). The bullets are
% the agreed framing; keep limitations quantitative and cite the
% oracle numbers, per the plan.

The oracle (Sec.~\ref{sec:results}) settles where the remaining error
lives. With perfect boxes the pipeline gains only \OracleDetection{}
terminal-pair $F_1$, and with perfect terminal snapping only
\OracleSnapping{}; injecting perfect \emph{connectivity} gains
\OracleWires{}. Symbol detection is solved on this dataset and terminal
snapping is close to exhausted; wire tracing is the bottleneck, and any
further work on this task should go there. This also illustrates a
methodological point worth stating: an earlier version of this analysis
attributed most of the error to snapping, and it was wrong for two
reasons — it predated the boundary-crossing snapper, and its synthetic
wire injection was subtly unreadable by the stage it was measuring.
Stage attribution is only as trustworthy as the verification of the
injected oracle, which is why we verify every synthetic render and
report how many were excluded.

\draftnote{remaining Discussion bullets to develop:
(i) why per-component accuracy is the number that gates strict success
--- with $\sim$14 components per image, strict success is a product
over components, so even a high per-component accuracy yields a modest
strict rate; state the arithmetic with generated numbers, not a
worked hypothetical;
(ii) errors are concentrated — rail-level fixes repair many components
at once, which is why stage fixes produce jumps, not creep;
(iii) what the ledger reveals: hand-drawn circuits are mostly
\emph{underspecified} rather than wrong — the common entries are
missing ground and unvalued parts, which is the argument for treating
intent completion as part of the problem;
(iv) stratified results — crossing-heavy and multi-terminal drawings
are measurably harder (Sec.~\ref{sec:results});
(v) comparison-by-citation with Wang/SINA: different, unreleased test
sets, so no direct comparison is possible — exactly the problem C1
fixes.}

\subsection*{Limitations}
\begin{itemize}
  \item \textbf{Single dataset; no drafter split.} Digitize-HCD has no
  drafter metadata (Sec.~\ref{sec:dataset}); results may not transfer
  across drawing styles. \draftnote{cite CGHD zero-shot if run.}
  \item \textbf{GT scope.} Connectivity GT covers the
  \NumTestImages-image test split, verified by one annotator
  \todoa{+ mentor second-read status}.
  \item \textbf{Ground-truth boxes are not geometric ground truth.} GT
  files were bootstrapped from pipeline output and verified for
  \emph{topology}; their bounding boxes were not. Roughly a fifth are
  square, which caps their IoU against a correctly-shaped detection and
  makes bounding-box alignment discard components the pipeline in fact
  got right. We report the alignment-threshold sensitivity rather than
  tuning the threshold to taste. \draftnote{decide before submission:
  re-project GT boxes from the published COCO annotations (the right
  fix, moves every number), or report the sensitivity curve and keep
  0.3. Do not do this silently either way.}
  \item \textbf{No value extraction.} Component values are
  placeholders, logged as assumptions (future work M4/OCR).
  \item \textbf{Repair assumptions are flagged, not validated.}
  Unless the expert-acceptance study runs, ledger entries are
  \emph{disclosed} choices, not human-endorsed ones — that is the
  honest claim.
  \item \textbf{Absolute strict success remains low} — this is a
  hard task and the benchmark is deliberately strict; the paper's
  claim is measured progress and a reusable yardstick, not a solved
  problem.
\end{itemize}
